\begin{figure*}[t]
    \makebox[\textwidth][c]{
    \begin{tikzpicture}
    \tikzstyle{every node}=[font=\scriptsize]
    \definecolor{tabblue}{RGB}{31, 119, 180}
\definecolor{taborange}{RGB}{255, 127, 14}
\definecolor{tabgreen}{RGB}{44, 160, 44}
\definecolor{tabred}{RGB}{214, 39, 40}
\definecolor{tabpurple}{RGB}{148, 103, 189}
\definecolor{tabbrown}{RGB}{140, 86, 75}
\definecolor{tabpink}{RGB}{227, 119, 194}
\definecolor{tabgray}{RGB}{127, 127, 127}
\definecolor{tabolive}{RGB}{188, 189, 34}
\definecolor{tabcyan}{RGB}{23, 190, 207}
\definecolor{lightblue}{RGB}{173, 216, 230}
\definecolor{sandybrown}{RGB}{244, 164, 96}
\definecolor{darkgrey}{RGB}{169, 169, 169}
\definecolor{dimgrey}{RGB}{105, 105, 105}
\definecolor{olivedrab}{RGB}{107, 142, 35}
\definecolor{darkviolet}{RGB}{148, 0, 211}
\definecolor{darkgoldenrod}{RGB}{184, 134, 11}
\definecolor{darkblue}{RGB}{0, 0, 139}
\definecolor{orchid}{RGB}{218, 112, 214}

    \begin{axis}[%
        hide axis,
        xmin=10,
        xmax=50,
        ymin=0,
        ymax=0.1,
        legend style={
            draw=white!15!black,
            legend cell align=left,
            legend columns=5,
            legend style={
                draw=none,
                column sep=1ex,
                line width=1pt,
            }
        },
        ]
        \addlegendimage{line legend, tabgreen, ultra thick} % Thicker line here
        \addlegendentry{\textbf{\model} (Ours)}
        \addlegendimage{line legend, sandybrown, ultra thick} % Thicker line here
        \addlegendentry{EMPN}
        \addlegendimage{line legend, dimgrey, ultra thick} % Thicker line here
        \addlegendentry{Transformer}
        \addlegendimage{line legend, darkblue, ultra thick} % Thicker line here
        \addlegendentry{HeteroGNN}
        \addlegendimage{line legend, orchid, ultra thick} % Thicker line here
        \addlegendentry{GNN}
    \end{axis}
\end{tikzpicture}

    }
    \centering
    % First row of figures
    \begin{subfigure}[b]{0.32\linewidth}
        \includegraphics[width=\textwidth]{ICLR_2025/Figures/eval_cloth_hanging_equi/eval_full_Isaac-Cloth-Hanging-Multi-v0_eval_all.pdf}
        \caption{Original 3D Space}
    \end{subfigure}
    \hfill
    \begin{subfigure}[b]{0.32\linewidth}
        \includegraphics[width=\textwidth]{ICLR_2025/Figures/eval_cloth_hanging_equi/eval_half_yaw_Isaac-Cloth-Hanging-Multi-v0_eval_all.pdf}
        \caption{$\theta_{\text{yaw}} \in (-\pi/2, \pi/2)$}
    \end{subfigure}
    \hfill
    \begin{subfigure}[b]{0.32\linewidth}
        \includegraphics[width=\textwidth]{ICLR_2025/Figures/eval_cloth_hanging_equi/eval_quater_yaw_Isaac-Cloth-Hanging-Multi-v0_eval_all.pdf}
        \caption{$\theta_{\text{yaw}} \in (-\pi/4, \pi/4)$}
    \end{subfigure}
    
    \medskip
    % Second row of figures - centered by wrapping in a minipage
    \begin{minipage}{0.65\textwidth}
    \centering
    \begin{subfigure}[b]{0.49\linewidth}
        \includegraphics[width=\textwidth]{ICLR_2025/Figures/eval_cloth_hanging_equi/appx_eval_quater_half_yaw_Isaac-Cloth-Hanging-Multi-v0_eval_all.pdf}
        \caption{$\theta_{\text{yaw}}\in (-\pi/8, \pi/8)$}
    \end{subfigure}
    \hfill
    \begin{subfigure}[b]{0.49\linewidth}
        \includegraphics[width=\textwidth]{ICLR_2025/Figures/eval_cloth_hanging_equi/appx_eval_one_yaw_Isaac-Cloth-Hanging-Multi-v0_eval_all.pdf}
        \caption{Fixed orientation}
    \end{subfigure}
    \end{minipage}
    \caption{Performance of different models on the \emph{Cloth-Hanging} task across various sample spaces. Assuming the global scene located at $r=[0,1,0]^T$, then from left to right, we generate sample by rotating $r$ by (a) $\theta_{\text{roll}} \in (-\pi/4, \pi/2)$, $\theta_{\text{yaw}} \in (-\pi, \pi)$, (b) $\theta_{\text{yaw}} \in (-\pi/2, \pi/2)$, and (c) $\theta_{\text{yaw}} \in (-\pi/4, \pi/4)$. Meanwhile, the bottom row shows results for (d) $\theta_{\text{yaw}}\in (-\pi/8, \pi/8)$, and (e) the fixed orientation at $\theta_{\text{roll}}=0, \theta_{\text{yaw}}=0$. As the sample space decreases, performance improves across all models, with HEPi consistently outperforming the baselines. The additional plot with fixed orientation on the bottom are averaged over 5 seeds while the others with 10 seeds.}
    \vspace{-0.2cm}
    \label{fig:appx_eval_equi}
\end{figure*}
